\chapter{Conclusion and Future Work}

\section{Conclusion}

This project built GAIA, a rover-assisted Generative AI system for proactive crop health monitoring, and evaluated it. The rover is
deliberately a thin data-acquisition layer; the engineering effort and the evaluation concentrate on two independent Generative AI
pipelines that share one knowledge store.

\textbf{Pipeline A} answers questions about the farmer's own field. It separates the roles of code and model: exact statistics, time
resolution and citation checks are done by deterministic code, while the model plans the query and phrases the answer. On a synthetic farm with ground truth
computed from the raw readings, the final configurations answered all ten evaluation questions with correct numbers, the right kind of evidence
and 100\,\% planner accuracy. \textbf{Pipeline B} turns a photograph into a possibility-phrased, evidence-based diagnosis with a single vision-language call and
one bounded research agent, and feeds the findings one way into the shared store, so that the assistant can later discuss a
diagnosis together with the conditions in the field.

The battery constraint of the rover led to a feature that matters beyond power saving: every sensor can be switched on or off,
the backend records the switches, readings and heartbeats, and the assistant explains gaps in the data from these records instead of guessing.
Streaming chat, a Live view of the pipeline, a rover simulator, bring-up tools and demonstration data complete the system, which is verified by
96 backend and 22 application automated tests and by both firmware sketches compiling without warnings.

Table~\ref{tab:objectives} relates the outcomes to the objectives of Section~1.5.

\begin{table}[H]
\centering
\caption{Objectives and outcomes.}
\label{tab:objectives}
\small
\begin{tabular}{|c|P{6.1cm}|P{5.2cm}|}
\hline
\textbf{Objective} & \textbf{Outcome} & \textbf{Evidence}\\\hline
O1 & Rover firmware, HTTP API, live video, stills & Compiled sketches; check tools\\\hline
O2 & Validated ingestion; windows, daily and on/off chunks; embeddings & Ingestion and PostgreSQL tests; Chapter~5\\\hline
O3 & Pipeline A with exact facts, hybrid retrieval, validated citations & 100\,\% on ten questions; Chapter~6\\\hline
O4 & Pipeline B with one vision call, one agent, one-way write & End-to-end tests and a local-model run\\\hline
O5 & Per-sensor switches and gap-aware answers & Switch tests; live simulator run\\\hline
O6 & Flutter application with streaming and Live view & Application tests; screenshots\\\hline
O7 & Responsible-AI practices & Prompts, validators, tests; Chapter~4\\\hline
O8 & Tests, evaluation and performance figures & Chapter~6 and Chapter~7\\\hline
\end{tabular}
\end{table}

\section{Limitations of the Proposed System}

\begin{itemize}
\item \textbf{Confidence is uncalibrated.} The model's self-reported confidence is only a rough indicator. The system says so on every result
      and provides the metrics to measure it, but it does not correct it.
\item \textbf{The vision model can be wrong.} It can misidentify a crop or a disease, especially on poor photographs; results are
      possibilities with evidence and look-alikes, not diagnoses.
\item \textbf{Search quality varies.} The agent reads search snippets, not full web pages, and results depend on the region and the language.
\item \textbf{Sensor limits.} The MQ sensors are uncalibrated and cross-sensitive and are reported as relative indices; the DHT11 has about $\pm$2\,$^{\circ}$C and
      $\pm$5\,\% RH accuracy; the soil probe needs per-probe calibration.
\item \textbf{No location.} Readings belong to the rover, not to a plot, because there is no GPS or odometry.
\item \textbf{Small evaluation.} The evaluation uses ten questions on a synthetic farm; the vision benchmark is an upper bound for field
      photographs, and results with the cloud model are reported only where they were measured.
\item \textbf{Power saving needs hardware.} A switch saves energy only when the sensor's supply passes through an electronic switch; otherwise the sensor is only ignored in software.
      Gas sensors need a few minutes to warm up after being switched on.
\item \textbf{Controller layout.} The firmware reads all sensors, including the ultrasonic one, on a single Nano; using a second Nano for the ultrasonic sensor would need an
      additional link to the ESP32.
\item \textbf{Single farm and user.} There are no accounts, and an optional shared API key is the only access control.
\item \textbf{Search latency.} The semantic search over the newest window lags by up to about 150 seconds after the window closes, although exact statistics are always current.
\end{itemize}

\section{Future Work}

\begin{itemize}
\item Calibrate the confidence (for example by mapping it to measured accuracy on a field-photograph set) and run a small usability study with farmers.
\item For uncertain results, ask for a better photograph before searching, which is currently an open design item.
\item Add a section in the application for farmers' own notes and documents (the backend already stores and retrieves them).
\item Let the chat offer a button that switches a sensor on when its tip is shown.
\item Track the rover's position (GPS or wheel odometry) so that readings and diagnoses belong to plots.
\item Add push alerts for unusual readings, adaptive sampling that raises the rate when conditions change, and solar charging.
\item Multi-tenant authentication, offline-first mobile synchronisation, and Indian-language voice interaction.
\item A larger evaluation: more questions, several cloud and local models, and a field-photograph benchmark.
\end{itemize}

\section{Applications}

\begin{itemize}
\item \textbf{Field scouting for small and medium farms.} Walk or drive the rover through a plot, measure conditions on demand and ask questions in plain language.
\item \textbf{Early warning of disease pressure.} Link high-humidity periods in the sensor log to photographs and diagnoses taken at the same time.
\item \textbf{Greenhouses and polyhouses.} The same pipeline can answer questions about temperature, humidity and gas levels in a closed space.
\item \textbf{Livestock sheds and biogas units.} The gas sensors can flag unusual readings near cattle sheds or digesters.
\item \textbf{Agricultural extension and education.} The system demonstrates grounded, honest generative AI, and the simulator and demonstration data run without any hardware.
\end{itemize}
